% !TeX root = ../../Book1.tex
% 纲要标题：### 6.4 中心扩张与群扩张的分裂性
% 纲要位置：计划纲要.md，第 296 行。

\section{中心扩张与群扩张的分裂性}
\label{b1:sec:0604}

同一个核和商群可以组成不同的扩张；即使核交换、商对核诱导的作用相同，中间群仍未必同构。循环群、二面体群与四元数群使这些差别可以直接计算：有的截面保持乘法，有的提升元素必须满足新的平方关系。中心扩张把共轭作用简化为恒等作用，因而尤其适合看清因子集所记录的额外信息。

% 本节正文按下列原纲要要点撰写。

% 纲要内容（仅作写作计划，尚非教材正文）：
%
% * 中心扩张与非分裂例子
% * 四元数群与二面体群的对照
% * 中心扩张的因子集与上同调解释
%
% ---
%

\subsection{中心扩张与非分裂例子}
\label{b1:subsec:0604:01}

\begin{definition}\label{b1:def:central-extension}
设 \(A,G,H\) 为群，给定同态 \(\iota:A\rightarrow G\) 和 \(\pi:G\rightarrow H\)，其中 \(\iota\) 单射、\(\pi\) 满射，且 \(\operatorname{im}\iota=\ker\pi\)。如果每个 \(\iota(a)\) 都与 \(G\) 中全部元素交换，即 \(\iota(A)\subseteq Z(G)\)，则称
\[
1\longrightarrow A\xrightarrow{\iota}G\xrightarrow{\pi}H\longrightarrow1
\]
为\term[zhongxin-kuozhang]{中心扩张}[central extension]。
\end{definition}

前三个条件说明这是一条短正合群列；新增加的中心性要求核中的元素逐个与全群交换，比核为正规子群更强。以下通过 \(\iota\) 把 \(A\) 看成 \(G\) 的子群。对任一归一化集合截面 \(s:H\rightarrow G\)，中心性使
\[
\alpha_h(a)=s(h)a\cdot s(h)^{-1}=a,
\]
所以 \(\alpha_h=\operatorname{id}_A\)。但截面代表的乘积仍可能与所选的乘积代表不同：
\[
s(h)s(k)=c(h,k)s(hk),\qquad c(h,k)\in A.
\]
中心性消去了共轭作用，却没有要求 \(c(h,k)=1\)；中心扩张因此仍可能不分裂。

\ProseParagraphBreak
中心扩张的核 \(A\) 必为交换群，因为 \(A\) 中元素属于 \(Z(G)\)，彼此交换；中间群 \(G\) 却不必交换。分裂时，中心性给出更强的结论。

\begin{proposition}[分裂中心扩张为直积]\label{b1:prop:split-central-extension}
设 \(1\rightarrow A\xrightarrow{\iota}G\xrightarrow{\pi}H\rightarrow1\) 为中心扩张，并通过 \(\iota\) 将 \(A\) 识别为 \(G\) 的子群。若该扩张分裂，则任一群同态截面 \(s:H\rightarrow G\) 都给出同构
\[
A\times H\longrightarrow G,\qquad(a,h)\longmapsto a\cdot s(h).
\]
因此，对于任意这样的中心扩张，若 \(H\) 交换而 \(G\) 非交换，则该扩张不分裂。
\end{proposition}

\begin{proof}
先假设扩张分裂。\cref{b1:thm:split-extension-semidirect}用于任一所给同态截面，给出 \(G\cong A\rtimes_\alpha H\)，其中 \(\alpha_h(a)=s(h)a\cdot s(h)^{-1}\)。因 \(A\leq Z(G)\)，有 \(\alpha_h(a)=a\)，半直积乘法遂成为 \((a,h)(a',h')=(aa',hh')\)，正是直积乘法。

若 \(H\) 也交换，则直积 \(A\times H\) 交换，所以 \(G\) 非交换时不可能分裂。
\end{proof}

这个同构保持指定扩张的结构：它把 \((a,1)\) 送到 \(\iota(a)\)，而 \(\pi(a\cdot s(h))=h\)。分裂判别要求同构与指定的核嵌入及商映射相容；仅给出一个抽象同构 \(G\cong A\times H\)，还没有核对这两项相容性。

\ProseParagraphBreak
即使 \(G\) 交换，扩张仍可能不分裂。例如\secref{b1:sec:0603}中的 \(C_4\rightarrow C_2\)：非零商元素的两个提升都为四阶，不能成为同态截面中二阶元素的像。

\ProseParagraphBreak
更一般地，对 \(m\geq2\) 考察加法群短正合列
\[
0\longrightarrow\mathbb Z\xrightarrow{\ \times m\ }\mathbb Z
\longrightarrow\mathbb Z/m\mathbb Z\longrightarrow0.
\]
这是中心扩张，但不分裂。若有同态截面，有限阶生成元 \([1]_m\) 的像 \(a\in\mathbb Z\) 必满足 \(ma=0\)，所以 \(a=0\)，不能映回 \([1]_m\)。障碍在于商群生成元满足有限阶关系 \(m[1]_m=0\)，而 \(\mathbb Z\) 中没有可满足 \(ma=0\) 的非零提升；这与 \(C_4\rightarrow C_2\) 中无法提升二阶关系是同一种现象。这里用 \(0\) 表示平凡加法群，与乘法群列两端的 \(1\) 含义相同。

\ProseParagraphBreak
核与商不决定中间群。例如核、商都为 \(C_2\) 时，\(C_2\times C_2\) 的坐标投影给出分裂中心扩张，\(C_4\) 的模 \(2\) 投影给出非分裂中心扩张。两个群可以由“非单位元素是否都平方为单位”区分，而这一差别没有包含在两个端点群的同构类型里。

\subsection{四元数群与二面体群}
\label{b1:subsec:0604:02}

下面先比较核为 \(C_4\)、商为 \(C_2\) 的两条扩张，检验核、商及商对核的诱导作用能否确定扩张。两者的诱导作用都是取逆，分裂性却不同；随后计算因子集，便能看到这些数据之外还缺少哪一条乘法关系。

\ProseParagraphBreak
四元数群 \(Q_8=\{\pm1,\pm i,\pm j,\pm k\}\) 满足
\(i^2=j^2=k^2=-1\)、\(ij=k=-ji\)。由\eqref{b1:eq:quaternion-relations}，\(i^2=-1\neq1\)、\(i^4=1\)，所以 \(\langle i\rangle\) 阶为 \(4\)；\cref{b1:prop:quaternion-subgroups}又说明它正规。其商有两个元素，故由\cref{b1:cor:finite-order-divides}同构于 \(C_2\)。以 \(C_4\rightarrow Q_8\) 将生成元映为 \(i\)，右端取商映射，得到
\[
1\longrightarrow C_4\longrightarrow Q_8\longrightarrow C_2\longrightarrow1.
\]
它不分裂：若有同态截面，\cref{b1:thm:split-extension-semidirect}给出一个映到商群 \(C_2\) 的同构补群，故核外必须有二阶元素。然而由 \(i^2=j^2=k^2=-1\) 可知六个元素 \(\pm i,\pm j,\pm k\) 都为四阶，唯一二阶元素 \(-1\) 已在核 \(\langle i\rangle\) 中，矛盾。

\ProseParagraphBreak
相对地，\(D_8\) 的旋转子群给出分裂扩张
\[
1\longrightarrow C_4\longrightarrow D_8\longrightarrow C_2\longrightarrow1,
\]
由\cref{b1:prop:dihedral-normal-form}，每个元素唯一写为 \(r^as^b\)，故 \(\langle r\rangle\cap\langle s\rangle=\{1\}\) 且 \(\langle r\rangle\langle s\rangle=D_8\)。商映射在反射子群 \(\langle s\rangle\) 上于是既单射又满射，故为同构，由\cref{b1:thm:split-extension-semidirect}给出分裂。

\ProseParagraphBreak
两个扩张的端点相同，商群 \(C_2\) 的非单位元素在核 \(C_4\) 上都诱导取逆自同构：二面体群有 \(srs^{-1}=r^{-1}\)，四元数群有 \(jij^{-1}=i^{-1}\)。因为核 \(C_4\) 交换，\cref{b1:cor:abelian-kernel-action}说明这个作用不随截面改变。因此，两个扩张具有相同的核、商及诱导作用，却有不同的分裂性。这里诱导的取逆作用确实给出群同态 \(C_2\rightarrow\operatorname{Aut}(C_4)\)，但 \(Q_8\) 的扩张仍不分裂；故截面所给的 \(\alpha\) 是群同态，并不足以推出该截面是同态或扩张分裂。

\ProseParagraphBreak
再看它们的中心。在 \(Q_8\) 中，\(\pm1\) 与全群交换，而六个元素 \(\pm i,\pm j,\pm k\) 各自都不与另一个基本元素交换，所以 \(Z(Q_8)=\{\pm1\}\)。在 \(D_8\) 中，\(r^a\) 与 \(s\) 交换当且仅当 \(a\) 为偶数，反射 \(r^as\) 则都不与 \(r\) 交换，故 \(Z(D_8)=\{1,r^2\}\)。两个中心商都同构于 \(C_2\times C_2\)，可直接给出同构：令 \((x,y)=(i,j)\) 或 \((r,s)\)，则
\[
C_2\times C_2\longrightarrow G/Z(G),\qquad
(a,b)\longmapsto x^ay^bZ(G).
\]
因为 \(x^2,y^2\in Z(G)\) 且 \(yx\) 与 \(xy\) 相差一个中心元素，该公式按模 \(2\) 指数良定义并保持乘法。\(x,y\) 生成 \(G\)，所以映射满射；两边都有四个元素，因此它是同构。因此又得到两条中心扩张
\[
1\longrightarrow C_2\longrightarrow G
\longrightarrow C_2\times C_2\longrightarrow1,
\qquad G=Q_8\ \text{或}\ D_8.
\]
这两条都不分裂：若分裂，\cref{b1:prop:split-central-extension}会给出 \(G\cong C_2\times(C_2\times C_2)\)，从而 \(G\) 交换；但 \(ij\neq ji\)、\(rs\neq sr\)，分别排除了 \(Q_8,D_8\) 的交换性。可见“这个群有一个分裂扩张表示”不等于“这个群的所有商映射都分裂”；\(D_8\) 的两条序列恰好说明这一点。

\ProseParagraphBreak
\cref{b1:tab:dihedral-quaternion-extensions}汇总这四条指定的扩张。各行的核都是交换群，所以可以谈论商群对核诱导的作用；后两行的核进一步包含于中心，作用才是平凡的。

\begin{table}[htbp]
\centering
\caption{二面体群与四元数群的四条扩张}
\label{b1:tab:dihedral-quaternion-extensions}
\begin{tabularx}{\linewidth}{@{}c>{\raggedright\arraybackslash}Xccc@{}}
\toprule
中间群 & 指定核 & 商群 & 诱导作用 & 分裂性 \\
\midrule
\(D_8\) & \(\langle r\rangle\cong C_4\) & \(C_2\) & 取逆 & 分裂 \\
\(Q_8\) & \(\langle i\rangle\cong C_4\) & \(C_2\) & 取逆 & 不分裂 \\
\(D_8\) & \(\{1,r^2\}\cong C_2\) & \(C_2\times C_2\) & 平凡 & 不分裂 \\
\(Q_8\) & \(\{1,-1\}\cong C_2\) & \(C_2\times C_2\) & 平凡 & 不分裂 \\
\bottomrule
\end{tabularx}
\end{table}

\hypertarget{b1:factor-comparison-deep}{因子集比较。}
先比较表中前两条扩张。把两条扩张的核都识别为 \(C_4=\langle z\rangle\)，其中 \(D_8\) 中令 \(z=r\)，\(Q_8\) 中令 \(z=i\)；把商群写成 \(C_2=\{1,q\}\)。因为 \(q^2=1\) 且归一化截面满足 \(\sigma(1)=1\)，有 \(c(q,q)=\sigma(q)^2\)：这个因子集值就是商生成元的所选提升的平方。在 \(D_8\) 中选 \(\sigma_D(q)=s\)，在 \(Q_8\) 中选 \(\sigma_Q(q)=j\)，则
\[
c_D(q,q)=s^2=1,
\qquad
c_Q(q,q)=j^2=-1=i^2=z^2.
\]
归一化条件使含单位元的因子集取值都为单位，所以 \(c(q,q)\) 是这里唯一可能非平凡的值；上述计算已经把分裂与不分裂两种平方关系写入因子集。

\ProseParagraphBreak
这两个值不能通过更换归一化截面而改变。两条扩张中 \(q\) 对 \(C_4\) 的作用 \(\alpha_q\) 都是取逆。若 \(\sigma'(q)=b(q)\sigma(q)\)，其中 \(b(q)\in C_4\) 且 \(b(1)=1\)，则由\eqref{b1:eq:factor-set-section-change}及 \(q^2=1\)，
\[
\begin{aligned}
c'(q,q)
&=b(q)\alpha_q\bigl(b(q)\bigr)c(q,q)b(1)^{-1}\\
&=b(q)b(q)^{-1}c(q,q)=c(q,q).
\end{aligned}
\]
因此 \(D_8\) 中的单位值保持为单位，而 \(Q_8\) 中的 \(z^2\) 不能靠改变截面消去。这说明两条扩张虽有相同的核、商和诱导作用，因子集仍记录了不同的分裂性。

\ProseParagraphBreak
再比较表中后两条中心扩张。此处重新把二阶中心核的生成元记为 \(z\)，选商群 \(C_2\times C_2\) 的两个生成元的提升 \(x,y\)，使
\[
z^2=1,\qquad yx=zxy,\qquad x^2=z,\qquad y^2=z^\delta,
\]
其中 \(D_8\) 取 \(x=r,y=s,\delta=0\)，\(Q_8\) 取 \(x=i,y=j,\delta=1\)。商群坐标记为 \((a,b)\)，其中 \(a,b\in\{0,1\}\)，选择截面 \(\sigma(a,b)=x^ay^b\)。计算 \(x^ay^bx^{a'}y^{b'}\) 时，额外的 \(z\) 来自三处：\(ba'\) 记录 \(y^b\) 移过 \(x^{a'}\) 的共轭关系；\(aa'\) 记录 \(x^2=z\) 产生的进位；\(\delta bb'\) 记录 \(y^2=z^\delta\) 产生的进位。把三项合并，得到
\[
\begin{aligned}
x^ay^b x^{a'}y^{b'}
&=z^{ba'}x^{a+a'}y^{b+b'}\\
&=z^{ba'+aa'+\delta bb'}
x^{(a+a')\bmod 2}y^{(b+b')\bmod 2},\\
c_\delta\bigl((a,b),(a',b')\bigr)
&=z^{ba'+aa'+\delta bb'}.
\end{aligned}
\]
指数模 \(2\)。由于商群的非零元素 \(h=(a,b)\) 都满足 \(h+h=(0,0)\)，截面提升的平方由因子集直接给出：
\[
\sigma(a,b)^2=c_\delta\bigl((a,b),(a,b)\bigr)
=z^{ab+a+\delta b}.
\]
核 \(\langle z\rangle\) 是二阶中心子群，\((z\sigma(h))^2=\sigma(h)^2\)，故同一个商元素的两个提升有相同的平方。三个非零商元素对应的结果为
\[
\begin{array}{c|cc}
 h & D_8\;\text{中提升的平方} & Q_8\;\text{中提升的平方}\\
 \hline
 (1,0) & z & z\\
 (0,1) & 1 & z\\
 (1,1) & 1 & z
\end{array}
\]
这个平方差同样不依赖归一化截面的选择。此处商群 \(C_2\times C_2\) 仍按坐标相加，其单位元为 \((0,0)\)，每个 \(h\) 都满足 \(h+h=(0,0)\)。把截面改为 \(\sigma'(h)=b(h)\sigma(h)\)，其中 \(b(h)\in\langle z\rangle\)、\(b(0,0)=1\)。中心扩张的诱导作用平凡，且二阶核中 \(b(h)^2=1\)，所以
\[
c'(h,h)=b(h)^2c(h,h)b(0,0)^{-1}=c(h,h).
\]
取 \(h=(0,1)\)，上式说明 \(D_8\) 中的 \(c_0(h,h)=1\) 与 \(Q_8\) 中的 \(c_1(h,h)=z\) 都不会随截面改变。若两条指定中心扩张等价，保持核与商识别的同构会把一条截面送为另一条扩张的截面，并保持这些因子集值，与上面的平方差矛盾。因此两条指定中心扩张不等价。

\ProseParagraphBreak
进一步，在 \(D_8\) 中，表内后两个商元素各有两个二阶提升，再加上中心的 \(z\)，共有五个二阶元素；在 \(Q_8\) 中，三个非零商元素的提升都为四阶，只有中心的 \(z\) 为二阶。同构保持元素阶，这才说明两个中间群作为抽象群也不同构。

\ProseParagraphBreak
作为对照，两群的共轭类大小分布却相同。在 \(Q_8\) 中，中心元素 \(1,-1\) 各自成类，其余三类为 \(\{\pm i\},\{\pm j\},\{\pm k\}\)，大小依次为 \(1,1,2,2,2\)。\secref{b1:sec:0402}已算出 \(D_8\) 也有这一分布。中心、中心商和共轭类大小分布都未区分它们，因子集给出的提升平方却能区分。

\subsection{中心扩张的因子集与上同调解释}
\label{b1:subsec:0604:03}

中心扩张中，所有 \(\alpha_h\) 都是恒等自同构。把作为交换群的核 \(A\) 改用加法记号，在\eqref{b1:eq:extension-factor-identity}中代入 \(\alpha_h=\operatorname{id}_A\)，因子集恒等式变为
\[
c(h,k)+c(hk,\ell)=c(k,\ell)+c(h,k\ell),
\]
并有 \(c(1,h)=c(h,1)=0\)。同样，把\eqref{b1:eq:factor-set-section-change}写为加法形式，更换归一化截面时有
\[
c'(h,k)=c(h,k)+b(h)+b(k)-b(hk).
\]
这里 \(b(1)=0\)。更换截面描述的是同一个扩张；比较不同扩张时，则按\cref{b1:def:extension-equivalence}要求中间群同构保持指定的核嵌入和商映射。下面的低阶计算可以同时说明这两种比较。

\ProseParagraphBreak
取 \(H=A=C_2\)，把两者都写为加法群 \(\{0,1\}\)，所有加法都按模 \(2\) 计算。此时归一化条件为 \(c(0,h)=c(h,0)=0\)，只有 \(c(1,1)=\varepsilon\) 可能非零。因子集恒等式在任一输入为 \(0\) 时自动成立；三个输入都为 \(1\) 时，两边同为 \(\varepsilon\)，所以 \(\varepsilon=0,1\) 都允许。按\cref{b1:prop:extension-factor-data}所给的形式，在 \(C_2\times C_2\) 的底集合上定义
\[
(a,h)(a',h')=\bigl(a+a'+c(h,h'),h+h'\bigr).
\]
将所得乘法集合记为 \(G_\varepsilon\)。还须验证它是群：三个元素相乘，两种结合次序的第一坐标分别多出
\(c(h,h')+c(h+h',h'')\) 与 \(c(h',h'')+c(h,h'+h'')\)，刚核验的因子集恒等式使它们相等，第二坐标本身满足结合律。归一化给出单位 \((0,0)\)；元素 \((a,h)\) 的逆为 \(\bigl(a+c(h,h),h\bigr)\)，左右相乘均为单位。投影到第二坐标是满同态，核为 \(\bigl\{(a,0):a\in C_2\bigr\}\)。由 \(c(0,h)=c(h,0)=0\)，核中的每个元素都与 \((a',h)\) 交换，所以得到所需中心扩张。

\ProseParagraphBreak
令 \(x=(0,1)\) 为非零商元素的所选提升，\(z=(1,0)\) 为核的非单位元素，则
\[
x^2=(\varepsilon,0).
\]
因此 \(\varepsilon\) 记录的正是提升平方。当 \(\varepsilon=0\) 时，\(x^2=(0,0)\)，乘法就是直积乘法，得到 \(G_0=C_2\times C_2\)。当 \(\varepsilon=1\) 时，\(x^2=z\)，四次方为 \((0,0)\)，故它生成整个四阶群，得到 \(G_1\cong C_4\)。对任意核与商均已识别为 \(C_2\) 的中心扩张，选一个归一化截面，由\cref{b1:prop:extension-factor-data}得到的坐标乘法必为以上两种之一；坐标对应保持核和商的识别，所以每条这样的扩张都等价于其中一条。

\ProseParagraphBreak
归一化截面改变只有 \(b(1)\) 一个参数，而
\(b(1)+b(1)-b(0)=0\)，不能改变 \(\varepsilon\)。也可以直接检验扩张等价的条件：若 \(F:G_\varepsilon\rightarrow G_{\varepsilon'}\) 保持核和商的识别，则 \(F(a,0)=(a,0)\)，且 \(F(0,1)=(u,1)\)，其中 \(u\in C_2\)。比较平方，得到
\[
F\bigl((0,1)^2\bigr)=(\varepsilon,0),\qquad
F(0,1)^2=(u+u+\varepsilon',0)=(\varepsilon',0).
\]
因此 \(\varepsilon=\varepsilon'\)，这两种模型给出恰好两个扩张等价类。其中 \(\varepsilon=0\) 的扩张分裂，而 \(\varepsilon=1\) 的因子集不能通过改变截面消去，扩张不分裂。

\ProseParagraphBreak
以下恢复一般群 \(H\) 的乘法记号，单位元写为 \(1_H\)；核 \(A\) 仍采用加法记号。第四卷第十二章“群的同调与上同调”将系统研究这种计算。固定交换群 \(A\)、群 \(H\) 及 \(H\) 在 \(A\) 上的平凡作用后，函数 \(c:H\times H\rightarrow A\) 满足上述归一化条件与结合律条件时，称为\term[guiyihua-er-yuquan]{归一化 \(2\)-余圈}[normalized 2-cocycle]。记
\[
Z^2(H,A)=\{\,c:H\times H\rightarrow A\mid
c\;\text{满足上述归一化与结合律条件}\,\}.
\]
\symidx{\(Z^2(H,A)\)：平凡作用下的归一化二次余圈群}
这些函数按逐点相加组成交换群：零函数是单位元，逐点取负给出逆元，而结合律条件对加法封闭。对任意满足 \(b(1_H)=0\) 的函数 \(b:H\rightarrow A\)，定义
\[
\delta b(h,k)=b(h)+b(k)-b(hk).
\]
\symidx{\(\delta b\)：由换截面产生的二次余边}
函数 \(\delta b\) 称为\term[er-yubian]{\(2\)-余边}[2-coboundary]；它正是改变归一化截面所产生的两个余圈之差。它满足归一化条件 \(\delta b(1_H,h)=\delta b(h,1_H)=0\)，且
\[
\begin{aligned}
\delta b(h,k)+\delta b(hk,\ell)
&=b(h)+b(k)+b(\ell)-b(hk\ell)\\
&=\delta b(k,\ell)+\delta b(h,k\ell).
\end{aligned}
\]
因此每个 \(2\)-余边都是归一化 \(2\)-余圈。映射 \(b\mapsto\delta b\) 保持逐点加法，所以
\[
B^2(H,A)=\{\,\delta b\mid b:H\rightarrow A,\ b(1_H)=0\,\}
\leq Z^2(H,A).
\]
\symidx{\(B^2(H,A)\)：平凡作用下的二次余边群}
\(Z^2(H,A)\) 收集所有满足结合律条件的归一化因子集，\(B^2(H,A)\) 收集仅由换截面产生的变化。把相差 \(2\)-余边的余圈视为同一类，得到第二上同调群
\[
H^2(H,A)=Z^2(H,A)/B^2(H,A).
\]
\symidx{\(H^2(H,A)\)：平凡作用下的第二上同调群}
这个商除去了截面选择造成的坐标差异。第四卷第12.2节将证明，它与固定 \(A,H\) 及其识别的中心扩张等价类一一对应：中间群的同构必须保持指定的核嵌入与商映射，满足\cref{b1:def:extension-equivalence}的条件。这里分类的是带这两条映射的扩张，而非单独按抽象同构比较中间群。
